\documentclass[11pt,oneside]{article}
\usepackage[a4paper,margin=27mm,headheight=15pt]{geometry}
\usepackage[T1]{fontenc}
\usepackage[utf8]{inputenc}
\IfFileExists{libertinus.sty}{\usepackage{libertinus}}{\usepackage{lmodern}}
\usepackage{microtype,amsmath,amssymb,amsthm,mathtools,mathrsfs}
\usepackage{booktabs,longtable,array,tabularx}
\usepackage{xcolor,fancyhdr,titlesec,enumitem,listings,xurl}
\usepackage{hyperref}
\usepackage[nameinlink,noabbrev]{cleveref}
\definecolor{linkblue}{RGB}{25,75,135}
\definecolor{shadegray}{RGB}{246,247,249}
\hypersetup{colorlinks=true,linkcolor=linkblue,citecolor=linkblue,urlcolor=linkblue,
 pdftitle={Recurrence-Free Dyadic Formulae for the Fabius and Rvachev Functions},
 pdfauthor={Companion research article prepared for Vladimir Reshetnikov},
 pdfsubject={Exact finite sums, Bernoulli partitions, binary telescopes, determinants, and quarter-base spline extraction}}
\pagestyle{fancy}\fancyhf{}
\fancyhead[L]{\small Recurrence-free dyadic formulae}
\fancyhead[R]{\small Fabius--Rvachev}
\fancyfoot[C]{\thepage}
\titleformat{\section}{\Large\bfseries}{\thesection}{0.7em}{}
\titleformat{\subsection}{\large\bfseries}{\thesubsection}{0.7em}{}
\setcounter{tocdepth}{2}
\setlength{\emergencystretch}{3em}
\setlist{itemsep=0.2em,topsep=0.4em}
\newtheorem{theorem}{Theorem}[section]
\newtheorem{lemma}[theorem]{Lemma}
\newtheorem{proposition}[theorem]{Proposition}
\newtheorem{corollary}[theorem]{Corollary}
\theoremstyle{definition}\newtheorem{definition}[theorem]{Definition}
\newtheorem{example}[theorem]{Example}
\theoremstyle{remark}\newtheorem{remark}[theorem]{Remark}
\newcommand{\R}{\mathbb R}\newcommand{\Q}{\mathbb Q}
\newcommand{\C}{\mathbb C}\newcommand{\N}{\mathbb N}
\newcommand{\Z}{\mathbb Z}\newcommand{\E}{\mathbb E}
\newcommand{\Prob}{\mathbb P}\newcommand{\Up}{\operatorname{up}}
\newcommand{\G}{\mathcal F}\newcommand{\eps}{\varepsilon}
\newcommand{\dd}{\,\mathrm d}\newcommand{\Comp}{\operatorname{Comp}}
\newcommand{\wt}{s_2}\newcommand{\D}{\mathrm D}
\newcommand{\qbinom}[3]{\genfrac{[}{]}{0pt}{}{#1}{#2}_{#3}}
\newcommand{\coef}[1]{[z^{#1}]}\newcommand{\Tn}{T_n}
\DeclareMathOperator{\supp}{supp}
\lstset{basicstyle=\ttfamily\small,columns=fullflexible,keepspaces=true,
 showstringspaces=false,breaklines=true,frame=single,backgroundcolor=\color{shadegray}}
\title{\bfseries Recurrence-Free Dyadic Formulae\\for the Fabius and Rvachev Functions\\[0.5em]
\large Explicit partitions, positive compositions, binary telescopes,\\integer determinants, and parity-reduced spline extraction}
\author{A companion research article prepared for Vladimir Reshetnikov}
\date{4 September 2026}
\begin{document}
\maketitle
\begin{abstract}
We give several completely finite descriptions of the exact rational values of
$F(a/2^n)$ and $\Up(a/2^n)$. Every auxiliary coefficient is specified by an explicit
finite sum or product; no moment recurrence or recursive evaluation of another
function value is required. The classical Thue--Morse formula is completed in two
ways: a Bernoulli partition expansion, with the Bernoulli numbers themselves
expanded into finite power sums, and a positive ordered-composition expansion.
A binary-digit telescope removes the long Thue--Morse prefix, while an integer
bordered determinant gives a compact arithmetic certificate and a common
denominator.

The main moment-free construction uses centered finite uniform splines. At a
dyadic argument of depth $n$, their cutoff dependence is an \emph{exact}
polynomial in $4^{-N}$ of degree at most $\lfloor n/2\rfloor$, not merely an
asymptotic expansion. Lagrange evaluation at zero therefore recovers the function
from only $\lfloor n/2\rfloor+1$ consecutive spline levels. Substitution of the
finite box-volume formula produces a double finite sum involving only integers,
factorials, Thue--Morse signs, and $q$-Pochhammer symbols at $q=1/4$.
We prove the cutoff-polynomial identity, its exact degree at reduced interior
dyadics, the connection with the existing half-base formula, and analogous
formulae for the Rvachev bump, derivatives, and iterated primitives. The same
construction extends to normalized reciprocal-integer-base uniform series.
Exact rational tests compare five formulae on 520 dyadic grid cases, together
with separate spline, bump, shift, and other-base tests. The accompanying
standard-library Python program reproduces these checks.
\end{abstract}
\clearpage
\begingroup\small
\tableofcontents
\endgroup
\newpage

\section{Scope, conventions, and the formula map}
\label{sec:scope}

\subsection{What is meant here by a closed form}
A finite expression is called \emph{recurrence-free} if its evaluation uses only
explicit finite index sets, finite sums and products, rational arithmetic,
integer powers, factorials, and directly specified combinatorial quantities.
Thus a formula involving moments $c_j$ is not yet complete until the $c_j$ have
been replaced by finite expressions. A coefficient extractor is useful in a
proof, but it is not left as an unexplained operation in the final evaluation
formulae. A determinant is permissible because we give its entries explicitly
and recall its finite permutation expansion.

Throughout, an empty sum is zero, an empty product is one, and $0^0=1$ in
polynomial formulae. All displayed index bounds are finite. We write
\begin{equation}
 \eps_k=(-1)^{\wt(k)},\qquad T_n=\binom{n+1}{2},
 \label{eq:sign}
\end{equation}
where $\wt(k)$ is the number of ones in the binary expansion of $k$.
This is a direct digit definition, not a recurrence. Equivalently, for any
$L$ with $k<2^L$,
\[
 \eps_k=\prod_{j=0}^{L-1}
 \left(1-2\left(\left\lfloor k/2^j\right\rfloor
                  -2\left\lfloor k/2^{j+1}\right\rfloor\right)\right).
\]
For $0<q<1$ we use the finite products
\begin{equation}
 (q;q)_r=\prod_{j=1}^{r}(1-q^j),\qquad
 \qbinom{d}{j}{q}=\frac{(q;q)_d}{(q;q)_j(q;q)_{d-j}}.
 \label{eq:qdef}
\end{equation}
No infinite $q$-product is needed to evaluate any dyadic value in this article.

\subsection{The three functions must be kept separate}
Let $U_1,U_2,\ldots$ be independent uniform random variables on $[0,1]$, and put
\begin{equation}
 X=\sum_{j=1}^{\infty}2^{-j}U_j,
 \qquad F(x)=\Prob(X\le x),\qquad Y=2X-1.
 \label{eq:random}
\end{equation}
The bounded CDF $F$ is zero on $(-\infty,0]$ and one on $[1,\infty)$.
The density of $Y$ is the even bump $\Up$, normalized by $\Up(0)=1$ and
$\supp\Up=[-1,1]$. In particular,
\begin{equation}
 \Up(t)=F(1-|t|)\quad(|t|\le1),
 \qquad F'(x)=2\Up(2x-1).
 \label{eq:dictionary}
\end{equation}
The signed extension is denoted by $\G$, not by $F$:
\begin{equation}
 \G(x)=\sum_{m\ge0}\eps_m\Up(x-2m-1).
 \label{eq:globaldef}
\end{equation}
This sum is locally finite, and at most one summand is nonzero. On $[0,1]$
it equals $F$, but on the rest of the real line it is not the bounded CDF.

Consequently, a formula for $F$ on the unit interval already answers the question
for every bump value:
\begin{equation}
 \boxed{\quad
 \Up(a/2^n)=
 \begin{cases}
 F\bigl((2^n-|a|)/2^n\bigr),& |a|\le2^n,\\
 0,&|a|>2^n.
 \end{cases}\quad}
 \label{eq:upreduction}
\end{equation}
For an arbitrary nonnegative dyadic $x$, set $m=\lfloor x/2\rfloor$ and
$t=x-2m$. Then
\begin{equation}
 \G(x)=\eps_m F(\min(t,2-t)),\qquad 0\le t<2.
 \label{eq:globalreduction}
\end{equation}
Negative arguments give $\G(x)=0$. These reductions avoid unnecessarily long
prefix sums outside the unit interval.

\subsection{Relation to the repository and earlier work}
The comparison source is the linked primary exposition in the \texttt{ProveIt}
repository, read at commit
\nolinkurl{2c6baf6738a5fbf3981a27141950edbe5952b03c}; its main TeX blob is
\nolinkurl{c5982d732a3caa769547e079efce37d25fb84a4e} \cite{repo}.
Its sections on reciprocal powers and arbitrary dyadics already contain moment
recurrences, a Thue--Morse expression, and exact path/composition descriptions
of $F(2^{-n})$. The archived companion \emph{The Dyadic Web of the Fabius
Function} already connects the half-base formula, binary reduction, moment
correction, and discrete extrapolation \cite{webcompanion}.

We therefore do not present the existence of composition formulae, the
Thue--Morse identity, or half-base interpolation as new discoveries.
The purpose here is to finish every coefficient into a finite expression,
organize genuinely different evaluation routes, and develop the centered,
parity-reduced cutoff identity with a precise level bound. Historical priority
for every consequence is not asserted. The classical analytic and arithmetic
background is developed by Arias de Reyna \cite{arias1982,arias2018};
Haugland gives another exact dyadic evaluator \cite{haugland}.
The half-base expression posed by Reshetnikov is discussed explicitly in
\cref{sec:halfbase} \cite{reshetnikov2019}.

\begin{table}[htbp]
\centering\small
\begin{tabularx}{\textwidth}{@{}p{0.23\textwidth}Xp{0.20\textwidth}@{}}
\toprule
Family & Finite ingredients and characteristic feature & Main result\\
\midrule
Bernoulli partitions & Multiplicity vectors and finite Bernoulli power sums; no recursive moments & \cref{thm:partitionvalue}\\
Positive compositions & Ordered positive parts and Mersenne factors; coefficient terms are positive & \cref{thm:compositionvalue}\\
Binary telescope & Only the nonzero binary digits of the argument; explicit moment coefficients & \cref{thm:binary}\\
Integer determinant & A matrix of order $\lfloor n/2\rfloor+1$; integer numerator certificate & \cref{thm:det}\\
Quarter-base extraction & $\lfloor n/2\rfloor+1$ centered spline levels; no moments at all & \cref{thm:quarter}\\
Direct bump extraction & The same weights, applied to finite densities instead of CDFs & \cref{thm:upquarter}\\
\bottomrule
\end{tabularx}
\caption{Alternative presentations. All rows are exact and recurrence-free.}
\label{tab:map}
\end{table}

\section{Analytic facts needed for exact arithmetic}
\label{sec:analytic}

\subsection{Smoothness, reflection, and the bump identity}
The series in \eqref{eq:random} converges absolutely for every outcome, because
$0\le U_j\le1$ and $\sum2^{-j}=1$. Replacing every $U_j$ by $1-U_j$ gives
$X\overset{d}=1-X$. Its law has no atoms: conditioning on all but $U_1$
leaves a translated uniform distribution. Hence
\begin{equation}
 F(1-x)=1-F(x)\qquad(0\le x\le1).
 \label{eq:reflection}
\end{equation}
The characteristic function of the centered variable is
\[
 \E e^{itY}=\prod_{j\ge1}\frac{\sin(t/2^j)}{t/2^j}.
\]
For any fixed $K$, the first $K$ factors give a bound
$C_K(1+|t|)^{-K}$; the remaining factors have absolute value at most one.
Choosing $K$ arbitrarily large makes the inverse Fourier integral and all its
derivatives absolutely convergent. Thus $Y$ has a $C^\infty$ density $\Up$,
supported in $[-1,1]$, and $F$ is smooth as well.

The self-similarity $Y\overset{d}=(V+Y')/2$, with $V$ uniform on $[-1,1]$
and $Y'$ an independent copy of $Y$, gives
\[
 \Up(t)=\Prob(2t-1\le Y'\le2t+1)=F(t+1)-F(t).
\]
On $[-1,0]$ this is $F(t+1)$; on $[0,1]$ it is $1-F(t)=F(1-t)$.
This proves \eqref{eq:dictionary}. The identity between the two densities
under $Y=2X-1$ gives $F'(x)=2\Up(2x-1)$. Differentiating the preceding
formula also gives
\begin{equation}
 \Up'(t)=2\Up(2t+1)-2\Up(2t-1).
 \label{eq:updiff}
\end{equation}
Smoothness and compact support imply that all positive-order derivatives of
$F$ vanish at $0$ and $1$, and all derivatives of $\Up$ vanish at $-1$ and $1$.

\subsection{Derivative termination at dyadic arguments}
\begin{lemma}[Signed extension and dyadic jets]
\label{lem:jets}
The function \eqref{eq:globaldef} is smooth and satisfies
\begin{equation}
 \G'(x)=2\G(2x),\qquad
 \G^{(r)}(x)=2^{r(r+1)/2}\G(2^r x).
 \label{eq:Gderiv}
\end{equation}
At nonnegative integers,
\begin{equation}
 \G(2m)=0,\qquad \G(2m+1)=\eps_m.
 \label{eq:Gintegers}
\end{equation}
If $0<a<2^n$, then
\begin{equation}
 F^{(r)}(a/2^n)=0\quad(r>n).
 \label{eq:Fjetvanish}
\end{equation}
Likewise, if $|a|<2^n$, then
\begin{equation}
 \Up^{(r)}(a/2^n)=0\quad(r>n).
 \label{eq:upjetvanish}
\end{equation}
\end{lemma}
\begin{proof}
Flatness at the endpoints of each translated bump makes the locally finite
sum smooth, including the joining points. By \eqref{eq:updiff}, its derivative
is
\[
 2\sum_{m\ge0}\eps_m\{\Up(2x-4m-1)-\Up(2x-4m-3)\}=2\G(2x),
\]
because the direct digit definition gives $\eps_{2m}=\eps_m$ and
$\eps_{2m+1}=-\eps_m$. Repeated differentiation proves
\eqref{eq:Gderiv}; evaluation at bump endpoints and centers gives
\eqref{eq:Gintegers}.
For an interior unit-interval argument, $F$ and $\G$ agree on a neighborhood.
When $r>n$, the number $2^r a/2^n$ is an even integer, so its signed value
vanishes. On $(-1,1)$, $\Up(x)=\G(x+1)$; the same reasoning applies to
$2^r(a/2^n+1)$. At $a=0,n=0$ this also proves that every positive-order
bump derivative vanishes at zero.
\end{proof}

\begin{remark}[A finite Taylor jet is not a local polynomial identity]
The vanishing in \eqref{eq:Fjetvanish} does not assert that $F$ equals its
Taylor polynomial near the dyadic point. The finite-spline construction below
uses a polynomial identity for a \emph{finite convolution}, together with an
exact smoothing identity. It does not replace $F$ by a convergent Taylor
series. This distinction is essential to the proof.
\end{remark}

\section{The master dyadic formula, proved from a finite cube}
\label{sec:master}

\subsection{A box-volume identity}
\begin{lemma}[Finite weighted-uniform CDF]
\label{lem:box}
For positive $w_1,\ldots,w_N$, let $Z=\sum_{j=1}^N w_jU_j$. Then
\begin{equation}
 \Prob(Z\le x)=\frac{1}{N!\prod_jw_j}
 \sum_{\delta\in\{0,1\}^N}(-1)^{|\delta|}
 \left(x-\sum_{j=1}^Nw_j\delta_j\right)_+^N.
 \label{eq:box}
\end{equation}
Here $s_+=\max(s,0)$ and $N\ge1$.
\end{lemma}
\begin{proof}
Changing variables $v_j=w_jU_j$ makes the desired probability the volume of
$\{0\le v_j\le w_j,\ \sum v_j\le x\}$ divided by $\prod_jw_j$.
The unrestricted nonnegative simplex has volume $x_+^N/N!$, as follows by
$N$ repeated integrations. Exclude the events $v_j>w_j$ by finite
inclusion--exclusion. For a selected set of excluded constraints, translate
the corresponding coordinates by $w_j$; the resulting simplex has threshold
$x-\sum w_j\delta_j$. Boundaries have measure zero. This gives
\eqref{eq:box}.
\end{proof}

Set
\begin{equation}
 X_N=\sum_{j=1}^N2^{-j}U_j,
 \qquad c_j=\E Y^{2j},\qquad C_j=\frac{c_j}{(2j)!},
 \qquad A_m=\frac{\E X^m}{m!}.
 \label{eq:moments}
\end{equation}
In particular, $C_0=A_0=1$. The symbols $C_j,A_m$ will shortly be replaced by
explicit sums.

\begin{theorem}[Master dyadic coefficient]
\label{thm:master}
For $n\ge1$ and $0\le a\le2^n$,
\begin{align}
 F(a/2^n)
 &=2^{-\binom n2}\sum_{h=0}^{a-1}\eps_h
       \sum_{m=0}^n A_m\frac{(a-h-1)^{n-m}}{(n-m)!},
 \label{eq:masterraw}\\
 &=2^{-T_n}\sum_{h=0}^{a-1}\eps_h
       \sum_{j=0}^{\lfloor n/2\rfloor}
             C_j\frac{(2a-2h-1)^{n-2j}}{(n-2j)!}.
 \label{eq:mastercentral}
\end{align}
Equivalently,
\begin{equation}
 F(a/2^n)=\frac{2^{-T_n}}{n!}
 \sum_{h=0}^{a-1}\eps_h\sum_{j=0}^{\lfloor n/2\rfloor}
   \binom n{2j}c_j(2a-2h-1)^{n-2j}.
 \label{eq:classical}
\end{equation}
For $n=0$, the same central expression gives the two allowed values $F(0)=0$
and $F(1)=1$.
\end{theorem}
\begin{proof}
Split the infinite series into independent head and tail:
$X\overset d=X_n+2^{-n}X'$. In \eqref{eq:box}, the weights $2^{-j}$ have
product $2^{-n(n+1)/2}$, and their subset sums are $h/2^n$ for
$0\le h<2^n$, with signs $\eps_h$. Conditioning on $X'$ therefore gives
\[
 F(a/2^n)=\frac{2^{-\binom n2}}{n!}
 \sum_{h=0}^{2^n-1}\eps_h\E(a-h-X')_+^n.
\]
If $h\ge a$, the summand is zero. If $h<a$, then $a-h\ge1$ and $0\le X'\le1$,
so the positive part can be removed. Reflection of $X'$ changes the
expectation into $\E(a-h-1+X)^n$. Expanding the ordinary binomial yields
\eqref{eq:masterraw}. Finally,
\[
 a-h-1+X=\tfrac12(2a-2h-1+Y),
\]
and odd moments of $Y$ vanish. The extra factor $2^{-n}$ converts
$2^{-\binom n2}$ into $2^{-T_n}$, proving the remaining formulae.
\end{proof}

This proof explains exactly why the tail produces only finitely many moments:
the dyadic threshold falls between consecutive head knots, so the positive
part becomes an ordinary polynomial throughout the whole tail support.
Rationality is not assumed anywhere in this argument.

\section{Removing all moment recurrences}
\label{sec:coefficients}

\subsection{Finite Bernoulli numbers and cumulants}
The standard Bernoulli convention is $B_1=-1/2$ \cite{dlmfBernoulli}.
To avoid hiding a recurrence even in this classical sequence, define
\begin{equation}
 B_m=\sum_{k=0}^m\frac{1}{k+1}
          \sum_{v=0}^k(-1)^v\binom kv v^m.
 \label{eq:Bernoullifinite}
\end{equation}
This finite expression agrees with the generating function
\begin{equation}
 \frac{z}{e^z-1}=\sum_{m\ge0}B_m\frac{z^m}{m!}.
 \label{eq:BernoulliGF}
\end{equation}
Indeed, the exponential generating function of the inner sum is
$(1-e^z)^k$. Summing $u^k/(k+1)=-\log(1-u)/u$ with $u=1-e^z$ gives
$z/(e^z-1)$. Coefficientwise, only $k\le m$ can contribute, since
$(1-e^z)^k$ has order $k$. Thus the argument is also valid as a formal
power-series identity, without an infinite numerical evaluation.

The moment generating functions are
\begin{equation}
 M_Y(z)=\E e^{zY}=\prod_{j\ge1}\frac{\sinh(z/2^j)}{z/2^j},
 \qquad M_X(z)=e^{z/2}M_Y(z/2).
 \label{eq:MGF}
\end{equation}
They are entire because the underlying variables are bounded. The logarithm
is needed only in a neighborhood of zero. Define the explicit rational numbers
\begin{equation}
 \lambda_r=\frac{4^r B_{2r}}{2r(2r)!(4^r-1)},\qquad
 \beta_r=\frac{B_{2r}}{2r(2r)!(4^r-1)}.
 \label{eq:lambdabeta}
\end{equation}
Then
\begin{equation}
 \log M_Y(z)=\sum_{r\ge1}\lambda_r z^{2r},
 \qquad
 \log M_X(z)=\frac z2+\sum_{r\ge1}\beta_r z^{2r}.
 \label{eq:logMGF}
\end{equation}
To verify the constants, differentiate $\log(\sinh z/z)$:
\[
 \coth z-\frac1z
 =\sum_{r\ge1}\frac{2^{2r}B_{2r}}{(2r)!}z^{2r-1}.
\]
Integrate with zero constant term and sum
$\sum_{j\ge1}2^{-2rj}=1/(4^r-1)$. This proves \eqref{eq:logMGF}.

\begin{theorem}[Partition formulae for the moments]
\label{thm:partitioncoeff}
For every $j,m\ge0$,
\begin{align}
 C_j&=\sum_{\substack{k_1,\ldots,k_j\ge0\\
                  k_1+2k_2+\cdots+jk_j=j}}
              \prod_{r=1}^j\frac{\lambda_r^{k_r}}{k_r!},
 \label{eq:Cpartition}\\
 A_m&=\sum_{\substack{k_0,k_1,\ldots,k_{\lfloor m/2\rfloor}\ge0\\
                  k_0+2\sum_{r\ge1}rk_r=m}}
       \frac{2^{-k_0}}{k_0!}
       \prod_{r=1}^{\lfloor m/2\rfloor}\frac{\beta_r^{k_r}}{k_r!}.
 \label{eq:Apartition}
\end{align}
At index zero, each sum consists of the unique empty or zero vector and
has value one.
\end{theorem}
\begin{proof}
Expand each exponential factor in
$M_Y(z)=\prod_{r\ge1}\exp(\lambda_rz^{2r})$ and collect the coefficient of
$z^{2j}$. Factors with $r>j$ cannot contribute. The second identity follows
in the same way from $e^{z/2}\prod_r\exp(\beta_rz^{2r})$. These are the
multiplicity-vector versions of the complete Bell-polynomial identities;
no recursively defined Bell polynomial is needed.
\end{proof}

\begin{theorem}[A fully expanded Bernoulli partition value]
\label{thm:partitionvalue}
Let $n\ge0$, $0\le a\le2^n$, and $d=\lfloor n/2\rfloor$. Then
\begin{equation}
 \boxed{\displaystyle
 F(a/2^n)=2^{-T_n}\sum_{h=0}^{a-1}\eps_h
 \sum_{\substack{k_0,k_1,\ldots,k_d\ge0\\
                 k_0+2\sum_{r=1}^d rk_r=n}}
 \frac{(2a-2h-1)^{k_0}}{k_0!}
 \prod_{r=1}^d\frac{\lambda_r^{k_r}}{k_r!}.}
 \label{eq:partitionvalue}
\end{equation}
Here every $\lambda_r$ is the rational expression \eqref{eq:lambdabeta},
with $B_{2r}$ explicitly given by \eqref{eq:Bernoullifinite}.
\end{theorem}
\begin{proof}
Substitute \eqref{eq:Cpartition} into \eqref{eq:mastercentral}, and set
$k_0=n-2j$. Conversely every vector in \eqref{eq:partitionvalue} determines
$j=\sum rk_r$, so the reindexing is bijective.
\end{proof}

\subsection{Positive ordered compositions: no Bernoulli numbers}
For $m\ge1$, let $\Comp(m)$ be the finite set of all ordered tuples
$\boldsymbol r=(r_1,\ldots,r_\ell)$ of positive integers with sum $m$.
Let $\Comp(0)$ contain only the empty tuple. Write
\begin{equation}
 S_i(\boldsymbol r)=r_i+\cdots+r_\ell.
 \label{eq:suffix}
\end{equation}
These suffix sums are part of the explicit tuple, not a recursively
constructed sequence. There are $2^{m-1}$ compositions of $m\ge1$, obtained
by choosing cuts among $m-1$ gaps.

\begin{theorem}[Positive composition coefficients]
\label{thm:compositioncoeff}
For $m,j\ge0$,
\begin{align}
 A_m&=\sum_{\boldsymbol r\in\Comp(m)}
       \prod_{i=1}^{\ell(\boldsymbol r)}
       \frac{1}{(r_i+1)!\,(2^{S_i(\boldsymbol r)}-1)},
 \label{eq:Acomp}\\
 C_j&=\sum_{\boldsymbol r\in\Comp(j)}
       \prod_{i=1}^{\ell(\boldsymbol r)}
       \frac{1}{(2r_i+1)!\,(4^{S_i(\boldsymbol r)}-1)}.
 \label{eq:Ccomp}
\end{align}
Every summand is positive.
\end{theorem}
\begin{proof}
First expand the finite products for $\E e^{zX_N}$. A nonconstant
contribution selects distinct levels $1\le j_1<\cdots<j_\ell\le N$ and
positive exponents $r_1,\ldots,r_\ell$ summing to $m$. Since
$\E U^r/r!=1/(r+1)!$, its coefficient is
\[
 \prod_{i=1}^\ell\frac{2^{-j_ir_i}}{(r_i+1)!}.
\]
The coefficients increase to $A_m$; equivalently, dominated convergence
applies to $X_N^m$. Put $j_0=0$ and $g_i=j_i-j_{i-1}\ge1$. Then
$\sum_i j_ir_i=\sum_i g_iS_i$, so the unrestricted level sum factors:
\begin{equation}
 \sum_{1\le j_1<\cdots<j_\ell}2^{-\sum_i j_ir_i}
 =\prod_{i=1}^{\ell}\sum_{g_i\ge1}2^{-g_iS_i}
 =\prod_{i=1}^{\ell}\frac{1}{2^{S_i}-1}.
 \label{eq:gaps}
\end{equation}
Only finitely many compositions remain, proving \eqref{eq:Acomp}.

For $Y=\sum_{j\ge1}2^{-j}V_j$ with $V_j$ uniform on $[-1,1]$, odd
coordinate moments vanish and
$\E V^{2r}/(2r)!=1/(2r+1)!$. Select positive even exponents $2r_i$.
The same gap calculation replaces $2^{S_i}$ by $4^{S_i}$ and yields
\eqref{eq:Ccomp}. One can first do the argument for a finite product and
then pass to the limit; absolute convergence and positivity of its even
coefficients justify the passage.
\end{proof}

\begin{theorem}[A Bernoulli-free closed dyadic sum]
\label{thm:compositionvalue}
For the same $a,n,d$ as above,
\begin{equation}
 \boxed{\begin{aligned}
 F(a/2^n)={}&2^{-T_n}\sum_{h=0}^{a-1}\eps_h
   \sum_{j=0}^{d}\frac{(2a-2h-1)^{n-2j}}{(n-2j)!}\\[-0.2em]
 &\hspace{1em}\times
 \sum_{\boldsymbol r\in\Comp(j)}
 \prod_{i=1}^{\ell(\boldsymbol r)}
 \frac{1}{(2r_i+1)!\,(4^{S_i(\boldsymbol r)}-1)}.
 \end{aligned}}
 \label{eq:compositionvalue}
\end{equation}
\end{theorem}
\begin{proof}
Insert \eqref{eq:Ccomp} into \eqref{eq:mastercentral}.
\end{proof}

Although the final prefix is signed, its universal coefficient calculation
is cancellation-free. In comparison with uncentered compositions, the largest
composition weight is only $\lfloor n/2\rfloor$, not $n$.
The first centered moments are
\begin{equation}
 c_0=1,\quad c_1=\frac19,\quad c_2=\frac{19}{675},\quad
 c_3=\frac{583}{59535},\quad c_4=\frac{132809}{32531625}.
 \label{eq:firstmoments}
\end{equation}
For example, the compositions $(2)$ and $(1,1)$ give
\[
 C_2=\frac{1}{5!(4^2-1)}
       +\frac{1}{3!\,3!\,(4^2-1)(4-1)}
 =\frac{19}{16200},
\]
and multiplication by $4!$ gives $c_2$.

\subsection{Reciprocal powers and two ways of eliminating the tail}
At $a=1$, \eqref{eq:masterraw} gives the particularly simple identity
\begin{equation}
 \boxed{\displaystyle F(2^{-n})=2^{-\binom n2}A_n.}
 \label{eq:inversepower}
\end{equation}
Thus either \eqref{eq:Apartition} or \eqref{eq:Acomp} is a direct finite
formula for these values. For odd depths, a second simplification follows
from the bump's half moments:
\begin{equation}
 F(2^{-(2j+1)})=
 \frac{c_j}{2(2j)!\,2^{\binom{2j+1}{2}}}.
 \label{eq:odddepth}
\end{equation}
To prove it without a moment recurrence, note that
$\Up(t)=\Prob(X\ge t)$ on $[0,1]$. Layer-cake integration gives
$\E X^{2j+1}=(2j+1)\int_0^1t^{2j}\Up(t)\dd t=(2j+1)c_j/2$.
Now apply \eqref{eq:inversepower}. This also explains why centered moments
can sometimes reduce the work more than raw moments do.

\section{A closed formula using only the nonzero binary digits}
\label{sec:binary}

The prefix in \eqref{eq:compositionvalue} has $a$ terms. For a large numerator
this is unnecessarily expensive: the binary expansion can replace it by at
most $n$ explicitly known polynomials. The mechanism is the exact Taylor
reduction used in the classical arithmetic theory \cite{arias2018}; here its
iteration is written as one finite sum and all its coefficients are eliminated.

\begin{lemma}[One binary step]
\label{lem:binarystep}
For an integer $b\ge1$ and $0\le t\le2^{-b}$,
\begin{equation}
 F(2^{-b}+t)+F(t)
 =2^{-\binom b2}\sum_{k=0}^{b}A_{b-k}\frac{(2^bt)^k}{k!}.
 \label{eq:binarystep}
\end{equation}
\end{lemma}
\begin{proof}
On $[0,2]$, the bump definition gives $\G(z+2)=-\G(z)$. Therefore the
$(b+1)$st derivative of the left side is zero: by \eqref{eq:Gderiv}, it
is a common constant times
$\G(2+2^{b+1}t)+\G(2^{b+1}t)$. The left side is consequently a polynomial
of degree at most $b$ on the stated interval. At $t=0$, every derivative of
$F(t)$ vanishes, whereas
\[
 F^{(k)}(2^{-b})
 =2^{k(k+1)/2}F(2^{-(b-k)})
 =2^{bk-\binom b2}A_{b-k}\qquad(0\le k\le b).
\]
The endpoint case $k=b$ uses $F(1)=A_0=1$. Its finite Taylor polynomial
is precisely \eqref{eq:binarystep}.
\end{proof}

\begin{theorem}[Explicit binary telescope]
\label{thm:binary}
Let $0<x<1$ be dyadic, and list its nonzero binary positions as
\[
 x=\sum_{i=1}^{r}2^{-b_i},\qquad 1\le b_1<\cdots<b_r.
\]
Define the explicitly known binary suffixes
\begin{equation}
 s_i=\sum_{j=i+1}^{r}2^{b_i-b_j},\qquad 0\le s_i<1.
 \label{eq:bitsuffix}
\end{equation}
Then
\begin{equation}
 \boxed{\displaystyle
 F(x)=\sum_{i=1}^{r}(-1)^{i-1}2^{-\binom{b_i}{2}}
       \sum_{k=0}^{b_i}\frac{s_i^k}{k!}
       \sum_{\boldsymbol r\in\Comp(b_i-k)}
       \prod_{j=1}^{\ell(\boldsymbol r)}
       \frac{1}{(r_j+1)!\,(2^{S_j(\boldsymbol r)}-1)}.}
 \label{eq:binaryclosed}
\end{equation}
The coefficient sum can equivalently be replaced by the finite partition
expression \eqref{eq:Apartition}.
\end{theorem}
\begin{proof}
Write $x_i=\sum_{j=i}^r2^{-b_j}$ and $x_{r+1}=0$. Since
$x_i=2^{-b_i}+x_{i+1}$ and $2^{b_i}x_{i+1}=s_i<1$, the one-step identity
says $F(x_i)+F(x_{i+1})=P_i(s_i)$, with $P_i$ the polynomial in
\eqref{eq:binarystep}. Multiply by $(-1)^{i-1}$ and sum over $i$.
All interior values cancel and $F(0)=0$. Substitution of
\eqref{eq:Acomp} gives the displayed formula. This is a proof by a finite
telescope, not an instruction to evaluate a recursive sequence.
\end{proof}

For example, $3/8=2^{-2}+2^{-3}$ gives
\[
 F(3/8)=2^{-1}\left(A_2+\tfrac12 A_1+\tfrac{1}{8}A_0\right)
        -2^{-3}A_3=\frac{73}{288}.
\]
After the explicit coefficients through degree $n$ have been computed once,
a point with $r$ nonzero bits uses $\sum_i(b_i+1)\le r(n+1)$ polynomial
terms. This count is not a bit-complexity estimate: large rational numerators
and coefficient generation still have costs. Reflection
$F(x)=1-F(1-x)$ may further reduce the number of useful bits.

\section{An integer determinant and an arithmetic certificate}
\label{sec:det}

A matrix formula packages the entire evaluation into one determinant with
integer entries. It is particularly convenient when a compact exact
certificate is preferable to a large expanded composition sum.

Put $d=\lfloor n/2\rfloor$ and define the integer power sums and row entries
\begin{equation}
 S_p(a)=\sum_{h=0}^{a-1}\eps_h(2a-2h-1)^p,
 \qquad t_j=\binom n{2j}S_{n-2j}(a)\quad(0\le j\le d).
 \label{eq:detrow}
\end{equation}
Let $L_d$ be the $d\times d$ lower triangular integer matrix with indices
$1\le k,j\le d$ and entries
\begin{equation}
 (L_d)_{kj}=
 \begin{cases}
 (2k+1)(4^k-1),&j=k,\\
 -\binom{2k+1}{2j},&j<k,\\
 0,&j>k.
 \end{cases}
 \label{eq:momentmatrix}
\end{equation}
Finally form the bordered matrix
\begin{equation}
 B_{n,a}=\begin{pmatrix}
 L_d&\boldsymbol 1\\
 -t_1\ \cdots\ -t_d&t_0
 \end{pmatrix},
 \label{eq:bordered}
\end{equation}
where $\boldsymbol 1$ has $d$ entries all equal to one. When $d=0$, this
means the one-by-one matrix $(t_0)$.

\begin{theorem}[Integer determinant formula]
\label{thm:det}
For $n\ge0$ and $0\le a\le2^n$,
\begin{equation}
 \boxed{\displaystyle
 F(a/2^n)=
 \frac{\det B_{n,a}}
 {2^{T_n}n!\displaystyle\prod_{k=1}^{d}(2k+1)(4^k-1)}.}
 \label{eq:detvalue}
\end{equation}
In particular, the reduced denominator of every value on this grid divides
\begin{equation}
 D_n=2^{T_n}n!(2d+1)!!\prod_{k=1}^{d}(4^k-1).
 \label{eq:denombound}
\end{equation}
No claim is made that $D_n$ is the smallest common denominator.
\end{theorem}
\begin{proof}
The centered MGF obeys
$M_Y(2z)=(\sinh z/z)M_Y(z)$, as follows either from the independent
uniform decomposition or directly by shifting the product factors.
Equating coefficients of $z^{2k}$, multiplying by $(2k+1)!$, and separating
the top term gives
\[
 (2k+1)(4^k-1)c_k=\sum_{j=0}^{k-1}\binom{2k+1}{2j}c_j.
\]
Since $c_0=1$, this is the finite system
$L_d(c_1,\ldots,c_d)^T=\boldsymbol 1$.
The diagonal entries are nonzero and
$\det L_d=\prod_{k=1}^d(2k+1)(4^k-1)$.
Subtracting a suitable linear combination of the first $d$ rows from the
last row of \eqref{eq:bordered}, or using a block determinant, gives
\[
 \det B_{n,a}=\det L_d\left(t_0+\sum_{j=1}^{d}t_jc_j\right).
\]
Now use \eqref{eq:classical}. The $d=0$ case follows directly from the
same formula. Every entry of $B_{n,a}$ is integral, proving the denominator
claim.
\end{proof}

The proof used a moment relation to identify a matrix; the \emph{final
expression} does not ask for any moments to be computed. In the literal
finite-sum interpretation,
\begin{equation}
 \det B_{n,a}=\sum_{\sigma\in S_{d+1}}
 (-1)^{\#\{(i,j):i<j,\ \sigma(i)>\sigma(j)\}}
 \prod_{i=1}^{d+1}(B_{n,a})_{i,\sigma(i)}.
 \label{eq:detfinite}
\end{equation}
Thus even the determinant notation can be removed. An implementation should
normally use exact elimination rather than this factorial-size expansion.
For $n=2,a=1$, the matrix is
\[
 B_{2,1}=\begin{pmatrix}9&1\\-1&1\end{pmatrix},
 \qquad F(1/4)=\frac{10}{2^3\,2!\,9}=\frac5{72}.
\]

\section{Centered splines and exact cutoff polynomials}
\label{sec:cutoff}

This section explains why a short moment-free formula exists. It is important
that the cutoff dependence terminates \emph{exactly}, not just to a prescribed
asymptotic accuracy.

\subsection{The finite centered CDF}
For $N\ge1$, set $h=2^{-N}$ and define
\begin{equation}
 \mathsf H_N(x)=\Prob(X_N+h/2\le x).
 \label{eq:centeredCDF}
\end{equation}
The finite random variable is centered at $1/2$, just like $X$.
Its CDF is a piecewise polynomial of degree $N$, with possible knots
$h(k+1/2)$, $0\le k<2^N$. For a dyadic argument $x=a/2^n$ and $N\ge n\ge1$,
write $M=a2^{N-n}$. The box formula immediately gives
\begin{equation}
 \boxed{\displaystyle
 \mathsf H_N(a/2^n)=
 \frac{1}{N!\,2^{N(N+1)/2}}
 \sum_{k=0}^{M-1}\eps_k(2M-2k-1)^N.}
 \label{eq:centeredCDFexplicit}
\end{equation}
The sum is empty at $a=0$. At $a=2^n$ it evaluates to one.
Although the density knots are numerous, the required threshold is always
halfway between two consecutive possible knots.

\subsection{A finite-jet deconvolution lemma}
The reciprocal centered MGF has the local series
\begin{equation}
 \frac{1}{M_Y(z)}=\sum_{j\ge0}R_jz^{2j},\qquad
 R_j=\sum_{\substack{k_1,\ldots,k_j\ge0\\\sum_{r=1}^jrk_r=j}}
      \prod_{r=1}^{j}\frac{(-\lambda_r)^{k_r}}{k_r!}.
 \label{eq:reciprocal}
\end{equation}
Here $R_0=1$, and $\lambda_r$ is the explicit number in
\eqref{eq:lambdabeta}. This coefficient formula follows by expanding
$\exp(-\sum_r\lambda_rz^{2r})$; it requires no reciprocal-series recurrence.

\begin{lemma}[Local polynomial smoothing is invertible on a finite jet]
\label{lem:localinverse}
Suppose a piecewise polynomial $H$ agrees with a polynomial $p(s)$ of degree
at most $N$ at the points $x+s$, $|s|\le\eta$. Let $Z$ be a symmetric random
variable supported on $[-1,1]$ with a smooth density, and suppose
$f(t)=\E H(t-\eta Z)$. Write
$M_Z(z)=\E e^{zZ}$ and $M_Z(z)^{-1}=\sum_{j\ge0}\rho_jz^{2j}$.
If the derivatives through order $N$ of the convolution are obtained by
convolving the corresponding piecewise derivatives of $H$, then
\begin{equation}
 H(x)=\sum_{j=0}^{\lfloor N/2\rfloor}
       \rho_j\eta^{2j}f^{(2j)}(x).
 \label{eq:localinverse}
\end{equation}
For the CDF splines used here, these derivative identities hold.
\end{lemma}
\begin{proof}
Let $\D$ denote differentiation of polynomials in $s$ and let
$J_Nf(s)=\sum_{r=0}^N f^{(r)}(x)s^r/r!$. Differentiating the convolution
as stipulated, and using the polynomial valid across the whole averaging
interval, shows coefficient by coefficient that
\[
 J_Nf=M_Z(-\eta\D)p=M_Z(\eta\D)p.
\]
This is an identity of polynomials, with the differential operator truncated
at order $N$. The second equality uses symmetry. Since $\D^{N+1}=0$ on
this space and $M_Z(0)=1$, its inverse is the finite operator obtained by
truncating $1/M_Z(\eta\D)$. Apply the inverse and evaluate at $s=0$.
This proves \eqref{eq:localinverse} without a Taylor-series convergence
assumption on $f$.

For an $N$-uniform CDF, derivatives through order $N-1$ are continuous;
the $N$th derivative is a bounded piecewise constant function. These are
also its distributional derivatives through order $N$. Convolution with a
smooth density therefore gives the derivative identities. Values at the
finitely many knot endpoints do not affect any integral. The same argument
works for a density spline, with its degree in place of $N$.
\end{proof}

\begin{theorem}[Exact polynomial in the cutoff scale]
\label{thm:cutoff}
Fix $n\ge1$, $0\le a\le2^n$, and $x=a/2^n$. For every $N\ge n$,
\begin{equation}
 \boxed{\displaystyle
 \mathsf H_N(x)=\sum_{j=0}^{d}
      R_j\,2^{-2j(N+1)}F^{(2j)}(x),
 \qquad d=\lfloor n/2\rfloor.}
 \label{eq:cutoffpoly}
\end{equation}
Thus there is one polynomial $P_x(z)$ of degree at most $d$, independent of
$N$, such that
\begin{equation}
 \mathsf H_N(x)=P_x(4^{-N}),\qquad P_x(0)=F(x).
 \label{eq:Pcutoff}
\end{equation}
\end{theorem}
\begin{proof}
The head--tail decomposition takes the centered form
\[
 X\overset d=X_N+h/2+(h/2)Y',\qquad
 F(t)=\E\mathsf H_N(t-(h/2)Y'),
\]
where $Y'$ is an independent copy of $Y$. Since $x=Mh$, the interval
$[x-h/2,x+h/2]$ contains no knot in its interior. On this entire interval
$\mathsf H_N(x+s)$ is represented by one polynomial of degree $N$;
continuity includes its endpoints. Apply Lemma~\ref{lem:localinverse} with
$\eta=h/2$ and $Z=Y'$. By Lemma~\ref{lem:jets}, all derivatives of $F$ at $x$
above order $n$ vanish. This removes all even terms with $2j>n$ and proves
\eqref{eq:cutoffpoly}. The endpoints $x=0,1$ can either be treated by the
same flatness statement or observed directly: both CDFs there are $0,1$.
The polynomial in \eqref{eq:Pcutoff} is explicitly
$P_x(z)=\sum_{j=0}^{d}R_j4^{-j}F^{(2j)}(x)z^j$.
\end{proof}

Only the constant coefficient of $P_x$ is needed. Its other coefficients
can be entirely eliminated by interpolation. For understanding the identity,
the first reciprocal coefficients are useful:
\begin{equation}
 R_0=1,\qquad R_1=-\frac1{18},\qquad
 R_2=\frac{31}{16200},\qquad R_3=-\frac{2347}{42865200}.
 \label{eq:firstR}
\end{equation}
For instance, $F''(1/4)=8$ and higher even derivatives vanish, giving the
particularly transparent exact formula
\begin{equation}
 \mathsf H_N(1/4)=\frac5{72}-\frac19\,4^{-N}\qquad(N\ge2).
 \label{eq:quarterexamplepoly}
\end{equation}

\subsection{The degree bound is sharp at reduced dyadic arguments}
\begin{proposition}[Nonzero reciprocal coefficients and exact degree]
\label{prop:exactdegree}
For every $j\ge0$, $(-1)^jR_j>0$. If $0<a<2^n$, $a$ is odd, and
$n\ge1$, the polynomial $P_{a/2^n}$ in \eqref{eq:Pcutoff} has degree
exactly $\lfloor n/2\rfloor$.
\end{proposition}
\begin{proof}
The classical product for $\sinh z/z$ \cite[\S4.36]{dlmfSinh} gives,
in a neighborhood of the origin,
\begin{equation}
 \frac1{M_Y(z)}=
 \prod_{\nu\ge1}\prod_{m\ge1}
 \left(1+\frac{z^2}{4^\nu\pi^2m^2}\right)^{-1}.
 \label{eq:reciprocalproduct}
\end{equation}
The positive numbers $\alpha_{\nu,m}=1/(4^\nu\pi^2m^2)$ have a finite
sum. Hence the absolute value of the coefficient of $z^{2j}$ is the convergent
strictly positive complete homogeneous sum of degree $j$ in this summable alphabet,
and its sign is $(-1)^j$. In particular it is nonzero. This infinite product
is used only to prove a sign; the evaluation formula for $R_j$ remains the
finite expression \eqref{eq:reciprocal}.

If $n=2d$, then
$F^{(2d)}(a/2^n)=2^{n(n+1)/2}\G(a)$ is nonzero because $a$ is odd.
If $n=2d+1$ and $d\ge1$, the relevant signed value is $\G(a/2)$, whose
magnitude is $F(1/2)=1/2$ by the bump decomposition. It too is nonzero.
For $n=1$, the constant polynomial has value $1/2$. Combining these facts
with $R_d\ne0$ proves the assertion.
\end{proof}

Reduction of the dyadic fraction is therefore useful: an unreduced depth
still gives a correct upper bound and a correct formula, but may introduce
unnecessary interpolation nodes.

\section{A moment-free quarter-base formula}
\label{sec:quarter}

\subsection{The geometric interpolation row}
For $0<q<1$ and $0\le j\le d$, define
\begin{align}
 w_{d,j}(q)
 &=\prod_{\substack{0\le\ell\le d\\\ell\ne j}}
        \frac{-q^\ell}{q^j-q^\ell}
 \label{eq:weightsproduct}\\
 &=\frac{(-1)^{d-j}q^{(d-j)(d-j+1)/2}}
        {(q;q)_j(q;q)_{d-j}}
 =\frac{(-1)^{d-j}q^{\binom{d-j+1}{2}}}{(q;q)_d}
       \qbinom{d}{j}{q}.
 \label{eq:weightsq}
\end{align}

\begin{lemma}[Constant extraction on a geometric grid]
\label{lem:weights}
For every polynomial $P$ of degree at most $d$ and every nonzero $z$,
\begin{equation}
 P(0)=\sum_{j=0}^{d}w_{d,j}(q)P(zq^j).
 \label{eq:geometricextract}
\end{equation}
Equivalently,
\begin{equation}
 \sum_{j=0}^{d}w_{d,j}(q)q^{jr}=\begin{cases}1,&r=0,\\0,&1\le r\le d.
 \end{cases}
 \label{eq:weightmoments}
\end{equation}
The generating polynomial of this row is
\begin{equation}
 \sum_{j=0}^{d}w_{d,j}(q)v^j
 =\prod_{r=1}^{d}\frac{v-q^r}{1-q^r}.
 \label{eq:weightpoly}
\end{equation}
\end{lemma}
\begin{proof}
The Lagrange basis polynomial for the node $zq^j$, evaluated at zero, is
exactly \eqref{eq:weightsproduct}; the common factor $z$ cancels. This
proves the extraction identity and its monomial specializations. Separating
the factors with $\ell<j$ and $\ell>j$ in the product gives
\eqref{eq:weightsq}. The generating polynomial has zeros at $q,\ldots,q^d$
by \eqref{eq:weightmoments}, and value one at $v=1$, proving
\eqref{eq:weightpoly}. These identities also follow from the finite
$q$-binomial theorem \cite[\S17.2]{dlmfq}, but no such theorem is needed
for this proof.
\end{proof}

\begin{theorem}[Quarter-base, moment-free dyadic evaluation]
\label{thm:quarter}
Let $n\ge1$, $0\le a\le2^n$, and $d=\lfloor n/2\rfloor$. Then
\begin{equation}
 \boxed{\displaystyle
 F(a/2^n)=\sum_{j=0}^{d}w_{d,j}(1/4)\,
                 \mathsf H_{n+j}(a/2^n).}
 \label{eq:quarterextract}
\end{equation}
Eliminating the splines gives the completely explicit double finite sum
\begin{equation}
 \boxed{\begin{aligned}
 F(a/2^n)={}&\sum_{j=0}^{d}
 \frac{(-1)^{d-j}4^{-\binom{d-j+1}{2}}}
 {(1/4;1/4)_j(1/4;1/4)_{d-j}\,(n+j)!\,2^{(n+j)(n+j+1)/2}}\\[-0.1em]
 &\hspace{1em}\times
 \sum_{k=0}^{2^j a-1}(-1)^{\wt(k)}
       (2^{j+1}a-2k-1)^{n+j}.
 \end{aligned}}
 \label{eq:quarterclosed}
\end{equation}
No moments, Bernoulli numbers, derivative values, recursively defined
coefficients, limits, or coefficient extractors occur in this formula.
\end{theorem}
\begin{proof}
Apply Lemma~\ref{lem:weights} to the degree-$d$ polynomial $P_x$ from
\cref{thm:cutoff}, with $z=4^{-n}$ and $q=1/4$.
Then substitute \eqref{eq:centeredCDFexplicit} at $N=n+j$.
The case $d=0$ uses the single weight one and needs no separate
interpolation argument.
\end{proof}

The first rows, useful for direct implementation, are
\begin{align*}
 d=0 &: \quad (1),\\
 d=1 &: \quad \frac13(-1,4),\\
 d=2 &: \quad \frac1{45}(1,-20,64),\\
 d=3 &: \quad \frac1{2835}(-1,84,-1344,4096).
\end{align*}
The largest spline level required is $n+\lfloor n/2\rfloor$.
The corresponding half-base construction has largest level $2n$.
This is a reduction in the number and size of finite splines, not a claim
that the method is faster than every other exact evaluator.

\subsection{Arbitrary starting depths and arbitrary nodes}
Consecutive levels need not start at $n$. For every $N_0\ge n$,
\begin{equation}
 F(a/2^n)=\sum_{j=0}^{d}w_{d,j}(1/4)\,
                  \mathsf H_{N_0+j}(a/2^n).
 \label{eq:quarterstart}
\end{equation}
More generally, for any distinct integers $N_0,\ldots,N_d\ge n$,
\begin{equation}
 F(a/2^n)=\sum_{j=0}^{d}
 \left(\prod_{\substack{0\le\ell\le d\\\ell\ne j}}
 \frac{-4^{-N_\ell}}{4^{-N_j}-4^{-N_\ell}}\right)
 \mathsf H_{N_j}(a/2^n).
 \label{eq:arbitrarycutoffs}
\end{equation}
These are immediate Lagrange evaluations of the same polynomial. They permit
reuse of already available spline values. The degree bound also explains the
universal node count: with $r\le d$ nonzero distinct nodes, the polynomial
$\prod_{i=1}^r(z-z_i)$ vanishes at every sampled node but not at zero.
Consequently no rule using only those samples can recover $P(0)$ for
\emph{every} polynomial of degree at most $d$. This is a statement about
universal polynomial recovery, not a lower bound for all special-function
algorithms.

\subsection{Size and signs of the interpolation coefficients}
The sign of $w_{d,j}(q)$ is $(-1)^{d-j}$. Substituting $v=-1$ into
\eqref{eq:weightpoly} therefore proves
\begin{equation}
 \sum_{j=0}^{d}|w_{d,j}(q)|
 =\prod_{r=1}^{d}\frac{1+q^r}{1-q^r}
 \le\prod_{r=1}^{\infty}\frac{1+q^r}{1-q^r}<\infty.
 \label{eq:weightnorm}
\end{equation}
The infinite product is finite because its logarithm is bounded by a
constant times $\sum_rq^r$. For $q=1/4$, the limiting amplification is
approximately $1.96926$, whereas for $q=1/2$ it is approximately $8.25599$.
Thus the interpolation step itself does not become progressively ill-conditioned
in the absolute maximum norm. Nevertheless, the alternating power sum inside
an individual spline can have severe cancellation. Exact integer and rational
arithmetic is the appropriate default. A nonnegative interpolation row is
impossible when $d\ge1$: positive nodes and nonnegative weights summing to
one cannot have zero first moment.

\section{How the half-base formula fits into the picture}
\label{sec:halfbase}

The quarter-base formula is not obtained by replacing $1/2$ by $1/4$ in an
existing identity. The centering removes all odd cutoff powers and changes
the required interpolation degree. A shift parameter makes the distinction
precise.

For $0\le\alpha\le1$, set
\begin{equation}
 \mathsf H_{N,\alpha}(x)=\Prob(X_N+\alpha2^{-N}\le x).
 \label{eq:shiftCDF}
\end{equation}
For $x=a/2^n$, $N\ge n\ge1$, and $M=a2^{N-n}$,
\begin{equation}
 \mathsf H_{N,\alpha}(a/2^n)
 =\frac{2^{-\binom N2}}{N!}
       \sum_{k=0}^{M-1}\eps_k(M-k-\alpha)^N.
 \label{eq:shiftCDFfinite}
\end{equation}
At $\alpha=0$ or $1$ the extra boundary terms have zero $N$th power,
so this same finite expression includes both endpoints. For arbitrary
complex $\alpha$, define $\mathsf H_{N,\alpha}$ at the displayed dyadic
argument by the polynomial on the right; it no longer has a probability
interpretation, but remains a finite algebraic object.

\begin{proposition}[General shift and half-base extraction]
\label{prop:halfbase}
For $n\ge1$, $0\le a\le2^n$, and every $\alpha\in\C$,
\begin{equation}
 F(a/2^n)=\sum_{j=0}^{n}w_{n,j}(1/2)\,
                   \mathsf H_{n+j,\alpha}(a/2^n).
 \label{eq:halfextract}
\end{equation}
Equivalently,
\begin{equation}
 \boxed{\begin{aligned}
 F(a/2^n)={}&\frac{1}{2^{n^2}(1/2;1/2)_n}
 \sum_{j=0}^{n}
 \frac{\qbinom{n}{j}{1/2}}{2^{j(j-1)}(n+j)!}\\[-0.1em]
 &\hspace{1em}\times
 \sum_{k=0}^{2^ja-1}\eps_k(k-2^ja+\alpha)^{n+j}.
 \end{aligned}}
 \label{eq:reshetnikov}
\end{equation}
\end{proposition}
\begin{proof}
First suppose $0\le\alpha\le1$ and put $h=2^{-N}$.
The head--tail identity reads
$F(t)=\E\mathsf H_{N,\alpha}(t-h(X'-\alpha))$.
At $x=Mh$, the entire averaging interval
$[x-h(1-\alpha),x+h\alpha]$ lies in one polynomial cell of
$\mathsf H_{N,\alpha}$. Its endpoints are consecutive possible knots.
Repeat the finite-jet proof of Lemma~\ref{lem:localinverse}, without symmetry.
Writing $Z=X'-\alpha$, the inverse operator is
$1/M_Z(-h\D)$, and therefore
\[
 \mathsf H_{N,\alpha}(x)=\sum_{r=0}^{n}
 \gamma_r(\alpha)h^rF^{(r)}(x),\qquad
 \gamma_r(\alpha)=[z^r]\frac1{M_Z(-z)}.
\]
Here $\gamma_0=1$ and the coefficients do not depend on $N$.
The formula is an exact polynomial in $2^{-N}$ of degree at most $n$.
Interpolation at $n+1$ consecutive depths proves \eqref{eq:halfextract}.
Both sides, when the finite formula \eqref{eq:shiftCDFfinite} is substituted,
are polynomials in $\alpha$. Equality throughout $[0,1]$ implies equality
for all complex $\alpha$.

Insert the explicit weights into \eqref{eq:halfextract} and replace
$(M-k-\alpha)^{n+j}$ by $(-1)^{n+j}(k-M+\alpha)^{n+j}$.
This cancels the weight's sign $(-1)^{n-j}$. The powers of two combine as
\[
 \binom{n-j+1}{2}+\binom{n+j}{2}=n^2+j(j-1),
\]
which gives \eqref{eq:reshetnikov}.
\end{proof}

Formula \eqref{eq:reshetnikov} is the unit-interval form of the expression
posed by Reshetnikov in 2019 \cite{reshetnikov2019}, subsequently developed
in the repository and its companion exposition \cite{repo,webcompanion}.
The repository also treats the signed extension directly. Our folding rule
\eqref{eq:globalreduction} is sufficient to evaluate it everywhere without
reproving that unrestricted-numerator identity.
At $\alpha=1/2$, $Z=Y'/2$ is symmetric and the odd coefficients vanish.
The degree drops from at most $n$ in $2^{-N}$ to at most
$\lfloor n/2\rfloor$ in $4^{-N}$. This is precisely the passage from
\eqref{eq:halfextract} to \eqref{eq:quarterextract}.

\section{Direct Rvachev formulae, derivatives, and primitives}
\label{sec:up}

\subsection{A finite density instead of a finite CDF}
Although \eqref{eq:upreduction} already determines every bump value, a direct
formula is useful when the available data are finite sinc products. Let
\begin{equation}
 Y_N=\sum_{j=1}^{N}2^{-j}V_j,\qquad V_j\sim\operatorname{Unif}[-1,1],
 \qquad u_N=\text{density of }Y_N.
 \label{eq:upspline}
\end{equation}
With the convention $\widehat f(t)=\int f(x)e^{-2\pi itx}\dd x$,
\begin{equation}
 \widehat u_N(t)=\prod_{j=0}^{N-1}
       \operatorname{sinc}(\pi t/2^j),\qquad
 \operatorname{sinc}(z)=\frac{\sin z}{z}.
 \label{eq:sincfinite}
\end{equation}
Thus an approximation indexed by the last factor $j=m$ is $u_{m+1}$ in
our notation. This offset matters in exact cutoff formulae.

For $|a|<2^n$ and $N>n\ge0$, put
\begin{equation}
 A=2^N+a2^{N-n}.
 \label{eq:upA}
\end{equation}
This is a positive even integer. The density version of the box formula is
\begin{equation}
 \boxed{\displaystyle
 u_N(a/2^n)=
 \frac{1}{(N-1)!\,2^{N(N-1)/2}}
 \sum_{k=0}^{A/2-1}\eps_k(A-1-2k)^{N-1}.}
 \label{eq:upsplineexplicit}
\end{equation}
Indeed, $Y_N=\sum_{j=1}^N2^{1-j}U_j-(1-2^{-N})$. Differentiate the
weighted-uniform CDF once; its subset thresholds are
$2^{-N}(2k-2^N+1)$. At the specified argument, the positive terms have
$2k<A-1$, giving precisely the displayed range and normalization.

\begin{theorem}[Direct quarter-base extraction for the bump]
\label{thm:upquarter}
Let $n\ge0$, $|a|<2^n$, and $d=\lfloor n/2\rfloor$. Then
\begin{equation}
 u_N(a/2^n)=\sum_{j=0}^{d}R_j4^{-Nj}\Up^{(2j)}(a/2^n)
                  \qquad(N>n),
 \label{eq:upcutoff}
\end{equation}
\begin{equation}
 \boxed{\displaystyle
 \Up(a/2^n)=\sum_{j=0}^{d}w_{d,j}(1/4)u_{n+1+j}(a/2^n).}
 \label{eq:upquarter}
\end{equation}
Each $u_{n+1+j}$ is the finite integer power sum
\eqref{eq:upsplineexplicit}. For $|a|\ge2^n$, the bump value is zero.
\end{theorem}
\begin{proof}
Now $Y\overset d=Y_N+hY'$ with $h=2^{-N}$, so
$\Up(t)=\E u_N(t-hY')$. The knots of $u_N$ are odd multiples of $h$.
Since $N>n$, the number $x/h=a2^{N-n}$ is even; hence
$[x-h,x+h]$ is one polynomial cell. On that cell $u_N$ has degree $N-1$.
Use the finite-jet inverse lemma with this degree and $\eta=h$.
All derivatives of $\Up$ above order $n$ vanish by Lemma~\ref{lem:jets},
giving \eqref{eq:upcutoff}. Constant extraction with $q=1/4$ yields
\eqref{eq:upquarter}. The case $N=1,n=0,a=0$ is the constant density on
its interior and follows by the same degree-zero argument. Outside
$[-1,1]$, and at its endpoints, the compactly supported smooth bump is zero.
\end{proof}

At $x=1/4$, the direct density construction yields
\[
 u_3(1/4)=\frac{15}{16},\qquad u_4(1/4)=\frac{179}{192},\qquad
 \Up(1/4)=\frac{4u_4(1/4)-u_3(1/4)}3=\frac{67}{72}.
\]
In fact $\Up''(1/4)=-8$, so the whole tail of the sequence satisfies
\begin{equation}
 u_N(1/4)=\frac{67}{72}+\frac49\,4^{-N}\qquad(N\ge3).
 \label{eq:upquarterexample}
\end{equation}
Unlike numerical extrapolation from an asymptotic expansion, two suitable
consecutive values determine the limit here with zero remainder.

\subsection{All derivatives at dyadics reduce to values}
The preceding formulae also give closed expressions for every derivative.
For an interior dyadic $x=a/2^n$ and $0\le r\le n$, use
\begin{align}
 F^{(r)}(x)&=2^{r(r+1)/2}\G(2^rx),
 \label{eq:Fderivativeclosed}\\
 \Up^{(r)}(x)&=2^{r(r+1)/2}\G(2^r(x+1)).
 \label{eq:upderivativeclosed}
\end{align}
In the first line, $x\in(0,1)$; in the second, $x\in(-1,1)$.
Fold the signed argument using \eqref{eq:globalreduction}, then evaluate
its unit-interval value by any finite formula above. The denominator depth
is at most $n-r$. All orders $r>n$ give zero. At $F$'s boundary points the
positive-order derivatives are zero; at the bump's boundary points all
orders, including zero, vanish.
This procedure consists of finite substitutions, not a chain of function-value
recurrences. It can be substituted directly into
\eqref{eq:cutoffpoly} or \eqref{eq:upcutoff} to make every polynomial
coefficient rational and explicit.

\subsection{Iterated primitives}
For an integer $r\ge1$, repeated integration has an equally simple exact
form:
\begin{align}
 \frac1{(r-1)!}\int_0^x(x-t)^{r-1}\G(t)\dd t
 &=2^{\binom r2}\G(x/2^r) &&(x\ge0),
 \label{eq:Gprimitive}\\
 \frac1{(r-1)!}\int_0^x(x-t)^{r-1}F(t)\dd t
 &=2^{\binom r2}F(x/2^r) &&(0\le x\le1),
 \label{eq:Fprimitive}\\
 \frac1{(r-1)!}\int_{-1}^x(x-t)^{r-1}\Up(t)\dd t
 &=2^{\binom r2}F((x+1)/2^r) &&(-1\le x\le1).
 \label{eq:upprimitive}
\end{align}
For \eqref{eq:Gprimitive}, differentiate the right side $r$ times and use
\eqref{eq:Gderiv}; the total power of two is
$\binom r2-r^2+r(r+1)/2=0$. Its first $r$ initial data at zero vanish,
so repeated FTC proves the identity. The second line restricts the first
to the unit interval. The third uses $\Up(t)=\G(t+1)$ and translates the
lower endpoint; the final argument belongs to $[0,1]$.
Thus these integrals at dyadic endpoints are rational and have exactly the
same recurrence-free evaluations at a larger denominator depth.

\section{The reciprocal-integer-base extension}
\label{sec:integerbase}

The same mechanism is not confined to base two. The dyadic case has the
extra convenience that all subset sums fill one consecutive integer block,
allowing a single Thue--Morse prefix. Other integer bases require a Boolean
cube sum instead.

\subsection{The normalized \texorpdfstring{base-$b$}{base-b} law}
Let $b\ge2$ be an integer and define
\begin{equation}
 X_b=(b-1)\sum_{j\ge1}b^{-j}U_j,\qquad
 F_b(x)=\Prob(X_b\le x).
 \label{eq:baselaw}
\end{equation}
It has support $[0,1]$, reflection symmetry, and a smooth density, by the
same absolutely convergent uniform-series and Fourier arguments used
above. In particular, it is flat at both boundary points. The raw moment
coefficients have the positive finite formula
\begin{equation}
 \boxed{\displaystyle
 A_m^{(b)}=\frac{\E X_b^m}{m!}
 =(b-1)^m\sum_{\boldsymbol r\in\Comp(m)}
       \prod_{i=1}^{\ell(\boldsymbol r)}
       \frac{1}{(r_i+1)!\,(b^{S_i(\boldsymbol r)}-1)}.}
 \label{eq:basecoeff}
\end{equation}
The proof repeats the gap summation \eqref{eq:gaps}, replacing $2$ by
$b$ and including the common scale $(b-1)^m$.

For $n\ge1$, $0\le a\le b^n$, and
$\delta=(\delta_0,\ldots,\delta_{n-1})\in\{0,1\}^n$, put
\begin{equation}
 \Delta_\delta=a-(b-1)\sum_{j=0}^{n-1}\delta_j b^j.
 \label{eq:basedelta}
\end{equation}

\begin{theorem}[Finite base-$b$ arithmetic]
\label{thm:basearithmetic}
For these parameters,
\begin{equation}
 \boxed{\displaystyle
 F_b(a/b^n)=\frac{b^{-\binom n2}}{(b-1)^n}
 \sum_{\substack{\delta\in\{0,1\}^n\\\Delta_\delta\ge1}}
 (-1)^{|\delta|}
 \sum_{m=0}^{n}A_m^{(b)}\frac{(\Delta_\delta-1)^{n-m}}{(n-m)!}.}
 \label{eq:basevalue}
\end{equation}
Together with \eqref{eq:basecoeff}, this is a finite rational expression.
Consequently $F_b$ takes rational values at every base-$b$ rational argument.
\end{theorem}
\begin{proof}
Split $X_b$ into its first $n$ coordinates and the independent tail
$b^{-n}X_b'$. Apply the finite box formula to the first coordinates,
whose weights have product $(b-1)^n b^{-n(n+1)/2}$. Conditional on the
tail, the truncated power is $b^{-n^2}(\Delta_\delta-X_b')_+^n$.
Since $\Delta_\delta$ is integral and $X_b'\in[0,1]$, it is zero for
$\Delta_\delta\le0$ and is an ordinary polynomial for
$\Delta_\delta\ge1$. Reflection replaces $1-X_b'$ by $X_b'$, so expansion
of $(\Delta_\delta-1+X_b')^n/n!$ gives the inner sum.
The scalar factors combine to the one displayed.
\end{proof}

\subsection{Parity-reduced extraction at \texorpdfstring{base $b^{-2}$}{base b to the power -2}}
For $N\ge1$, set
\[
 X_{b,N}=(b-1)\sum_{j=1}^N b^{-j}U_j,
 \qquad \mathsf H_{N,b}(x)=\Prob(X_{b,N}+b^{-N}/2\le x).
\]
For $N\ge n$, let $M=ab^{N-n}$. Its direct finite evaluation is
\begin{equation}
 \begin{split}
 \mathsf H_{N,b}(a/b^n)
 ={}&\frac{1}{2^N(b-1)^N b^{\binom N2}N!}\\
 &\times\sum_{\delta\in\{0,1\}^N}(-1)^{|\delta|}
 \left(2M-1-2(b-1)\sum_{j=0}^{N-1}\delta_jb^j\right)_+^N.
 \end{split}
 \label{eq:basespline}
\end{equation}
This is still a finite rational formula; positive part here is just a
comparison of explicitly given integers.

\begin{theorem}[Moment-free integer-base extraction]
\label{thm:basequarter}
For $n\ge1$, $0\le a\le b^n$, and $d=\lfloor n/2\rfloor$,
\begin{equation}
 \boxed{\displaystyle
 F_b(a/b^n)=\sum_{j=0}^{d}w_{d,j}(b^{-2})\,
                       \mathsf H_{n+j,b}(a/b^n).}
 \label{eq:basequarter}
\end{equation}
All quantities on the right are finite products or the Boolean cube sum
\eqref{eq:basespline}.
\end{theorem}
\begin{proof}
The self-similarity
$X_b\overset d=((b-1)U+X_b')/b$ gives
\begin{equation}
 F_b'(x)=\frac b{b-1}\{F_b(bx)-F_b(bx-(b-1))\}.
 \label{eq:basediff}
\end{equation}
Iterating this differentiation identity $n$ times writes $F_b^{(n)}(x)$
as a finite linear combination of values $F_b(b^nx-k)$ with integral $k$.
Further differentiation and evaluation at $x=a/b^n$ involves only
positive-order derivatives of $F_b$ at integers. All such derivatives
vanish: integers are boundary points or lie outside the support of the
density. Thus $F_b^{(r)}(a/b^n)=0$ for $r>n$.

Write $Z_b=2X_b-1$, a symmetric variable supported on $[-1,1]$, and
$h=b^{-N}$. The head--tail identity is
$F_b(t)=\E\mathsf H_{N,b}(t-(h/2)Z_b')$.
Every knot of the centered finite CDF is an odd half-integer multiple of
$h$, while $a/b^n$ is an integer multiple of $h$ for $N\ge n$.
The averaging interval therefore lies in one polynomial cell; missing
knots cause no problem. The finite-jet inverse lemma proves that
$\mathsf H_{N,b}(a/b^n)$ is an exact polynomial in $b^{-2N}$ of degree
at most $d$, with constant coefficient $F_b(a/b^n)$. Apply geometric
constant extraction with $q=b^{-2}$.
\end{proof}

For example, the explicit sums give
\[
 F_3(1/9)=\frac7{576},\quad F_3(1/27)=\frac1{6912},\quad
 F_4(1/16)=\frac1{240},\quad F_4(1/64)=\frac1{51840}.
\]
At $b=2$, the cube sums collapse to the dyadic prefix and all preceding
formulae are recovered. For $b>2$, some of the highest jets can vanish
for additional geometric reasons; we assert the degree \emph{upper bound},
not the dyadic exact-degree statement. In particular, this extension does
not assert that the base-$b$ functions have all the same local analytic
properties as the dyadic Fabius function.

\section{Worked exact values and reproducible verification}
\label{sec:tests}

\subsection{Small closed-form examples}
The first quarter-base row already gives
\begin{equation}
 F(1/4)=\frac{4(13/192)-1/16}{3}=\frac5{72},
 \label{eq:examplefourth}
\end{equation}
where $\mathsf H_2(1/4)=1/16$ and $\mathsf H_3(1/4)=13/192$.
For a denominator of eight, the same two weights apply:
\[
 \mathsf H_3(3/8)=\frac{97}{384},\qquad
 \mathsf H_4(3/8)=\frac{389}{1536},\qquad
 F(3/8)=\frac{4\mathsf H_4(3/8)-\mathsf H_3(3/8)}3
       =\frac{73}{288}.
\]
At depth four, three levels suffice:
\[
 \mathsf H_4(1/16)=\frac1{24576},\qquad
 \mathsf H_5(1/16)=\frac{121}{1966080},\qquad
 \mathsf H_6(1/16)=\frac{6331}{94371840}.
\]
\[
 F(1/16)=\frac{\mathsf H_4(1/16)-20\mathsf H_5(1/16)
                         +64\mathsf H_6(1/16)}{45}
       =\frac{143}{2073600}.
\]
Each spline fraction comes from a finite signed sum of odd integer powers;
no previously known value of $F$ was inserted into the extraction.

\begin{table}[htbp]
\centering
\begin{tabular}{@{}rr@{\qquad}rr@{}}
\toprule
$a$ & $F(a/16)$ & $a$ & $F(a/16)$\\
\midrule
1 & $143/2073600$ & 5 & $305857/2073600$\\
2 & $1/288$ & 6 & $73/288$\\
3 & $46657/2073600$ & 7 & $777743/2073600$\\
4 & $5/72$ & 8 & $1/2$\\
\bottomrule
\end{tabular}
\caption{Exact dyadic values. The remaining interior entries follow by
$F((16-a)/16)=1-F(a/16)$. The bump values follow from
$\Up(a/16)=F((16-|a|)/16)$.}
\label{tab:sixteenths}
\end{table}

\subsection{What the accompanying program verifies}
The file \texttt{verify.py} uses only the Python standard library and
\texttt{fractions.Fraction}. It implements five separate evaluations:
the centered composition formula, the centered partition formula, the
binary telescope, the bordered determinant, and the quarter-base spline
formula. The determinant is evaluated by exact rational elimination.
Composition sets are enumerated by cut masks, and partition multiplicities
by a bounded Cartesian product; neither coefficient generator uses a
moment recurrence. Memoization merely avoids repeating an explicit sum.

The completed verification run recorded in \texttt{verification.json}
contains the following counts.
\begin{center}
\small
\begin{tabularx}{\textwidth}{@{}l r X@{}}
\toprule
Check & Cases & Range or description\\
\midrule
Centered coefficient equality & 9 & $C_j$ from compositions and partitions,
$0\le j\le8$\\
Five dyadic value formulae & 520 & Every $0\le a\le2^n$, $0\le n\le8$;
reflection and refinement are checked too\\
Exact cutoff polynomial & 360 & Interior points at $1\le n\le6$,
with $N=n,n+1,n+3$\\
Direct bump extraction & 132 & Every $-2^n\le a\le2^n$, $0\le n\le5$\\
Half-base shift independence & 65 & Selected points through depth five,
$\alpha=0,1/2,1,-2,2/3$\\
Other integer bases & 1252 & Every $0\le a\le b^n$ for $b=3,4,5$ and
$1\le n\le4$\\
\bottomrule
\end{tabularx}
\end{center}
The 520 grid cases are numerator--depth pairs, not 520 distinct rational
arguments: unreduced representations are intentionally included.
Eight standard reciprocal-power values are checked separately as anchors.
All comparisons are exact equality of rational numbers; there is no
floating-point tolerance. These tests check the implementations and indexing.
They are not a replacement for the proofs, nor are they a claim that the new
presentation has been formalized in Lean.

To reproduce the report, run
\begin{lstlisting}[language=bash]
python verify.py --json verification.json
\end{lstlisting}
The default maximum dyadic depth is eight. A larger value can be requested
with \texttt{--max-depth}; the full-grid tests grow exponentially.
The other test ranges remain fixed.

\subsection{A minimal quarter-base evaluator}
The following self-contained version shows how little algebra is required
for the moment-free formula. The full verification program supplies the
other methods and cross-checks.
\begin{lstlisting}[language=Python,basicstyle=\ttfamily\footnotesize]
from fractions import Fraction
from math import factorial

def fabius_quarter(a: int, n: int) -> Fraction:
    """Exact bounded Fabius value at a / 2**n."""
    if not isinstance(a, int) or not isinstance(n, int):
        raise TypeError("a and n must be integers")
    if n < 0 or not 0 <= a <= 2**n:
        raise ValueError("require n >= 0 and 0 <= a <= 2**n")
    if a == 0 or a == 2**n:
        return Fraction(int(a != 0))
    # Reduce the denominator before choosing the number of levels.
    while n > 0 and a % 2 == 0:
        a //= 2
        n -= 1
    d = n // 2
    q = Fraction(1, 4)
    result = Fraction(0)
    for j in range(d + 1):
        weight = Fraction(1)
        for ell in range(d + 1):
            if ell != j:
                weight *= -q**ell / (q**j - q**ell)
        N = n + j
        M = a * 2**j
        numerator = sum(
            (-1 if bin(k).count("1") % 2 else 1)
            * (2*M - 2*k - 1)**N
            for k in range(M)
        )
        spline = Fraction(numerator,
                          factorial(N) * 2**(N*(N+1)//2))
        result += weight * spline
    return result
\end{lstlisting}
This evaluator makes no call to a function-value or moment recurrence.
The finite loops implement the two sums and the weight product literally.
Reduction of the input is optional mathematically, but useful computationally.

\section{Choosing a representation and possible extensions}
\label{sec:conclusions}

\subsection{Which finite formula to use}
For a single value with a sparse binary expansion, the binary telescope
usually avoids the largest prefix sums. For many values on one grid, the
centered moment formula permits the same explicitly computed $C_j$ to be
reused. The composition expression is particularly transparent when positive
coefficient terms and elementary denominators are desirable. The partition
expression connects directly to cumulants and is well suited to symbolic
manipulation. The determinant is a compact integer certificate, with a
visible common denominator. The quarter-base construction is the most direct
answer when \emph{no auxiliary arithmetic coefficients at all} are desired:
it contains only a short geometric interpolation row and finite power sums.

These are algebraically equivalent descriptions, but they are not identical
algorithms. The number of summands does not alone measure performance.
Large powers and cancellation, bit lengths of rational intermediates,
coefficient reuse, and whether one value or a whole grid is required all
matter. The exact arithmetic in the program is intentionally straightforward;
it is a verification reference rather than a claim of optimal implementation.

\subsection{Further deductions that do not require additional conjectures}
A finite subset of the available cutoff coefficients can vanish at special
points. Whenever the nonzero orders are known, one may interpolate only
those monomials rather than the complete degree space. The exact-degree
result at reduced dyadics shows that the highest allowed order does not
vanish, but does not exclude gaps at lower orders. This suggests
point-adapted extraction rows, with their validity proved by the corresponding
finite moment equations.

The local inverse proof applies more generally whenever a smooth law is an
independent sum of a finite spline head and a scaled tail, the support of the
tail fits in one polynomial cell at the evaluation point, and the relevant
jet of the limiting function terminates. Symmetry removes the odd cutoff
powers; other coefficient zeros can remove further powers. One can therefore
separate the search for additional exact formulae into a geometric question
about knot cells and an algebraic question about the nonzero jet orders.
The integer-base theorem is one concrete realization of this separation.

\subsection{Conclusion}
The requested rational values admit several closed forms in the strict
finite sense. The hidden moment recurrences in the classical formula can
be eliminated by explicit multiplicity partitions, by positive ordered
compositions, or by an integer determinant. Binary digits provide a finite
telescope with no long signed prefix. Most importantly, centered finite
convolutions have an exact polynomial dependence on the squared tail scale
at dyadic arguments. Removing that polynomial by geometric interpolation
gives \eqref{eq:quarterclosed}, a moment-free $q=1/4$ formula requiring
only $\lfloor n/2\rfloor+1$ spline levels. Reflection and folding cover all
Rvachev and signed Fabius values, and the same finite algebra extends to
all derivatives, normalized iterated primitives, and reciprocal-integer-base
uniform laws.

\appendix
\clearpage
\section{A direct finite definition of every index set}
\label{app:indices}
To remove any ambiguity about the word ``finite'', all auxiliary indexing
can be reduced to bounded integer tuples. A composition of $m\ge1$ is
specified by a length $1\le\ell\le m$ and a tuple
$(r_1,\ldots,r_\ell)\in\{1,\ldots,m\}^{\ell}$ with
$r_1+\cdots+r_\ell=m$. The empty tuple is the only composition of zero.
A partition multiplicity vector of weight $j$ lies in
\[
 \prod_{r=1}^{j}\{0,1,\ldots,\lfloor j/r\rfloor\}
 \quad\text{and satisfies}\quad \sum_{r=1}^{j}rk_r=j.
\]
In a raw partition of weight $m$, add
$k_0\in\{0,\ldots,m\}$ and impose $k_0+2\sum_r rk_r=m$.
All binomial coefficients can be read as factorial quotients in their
stated ranges. All Bernoulli numbers are the finite power sums
\eqref{eq:Bernoullifinite}. Every binary sign is the explicit digit product
following \eqref{eq:sign}. Every determinant is the permutation sum
\eqref{eq:detfinite}. Every $q$-Pochhammer symbol in a final value formula
is the finite product \eqref{eq:qdef}. There is therefore no unstated
recursive sequence anywhere in the evaluation layer.

\begingroup
\small
\begin{thebibliography}{99}
\setlength{\itemsep}{3pt}
\bibitem{repo}
Vladimir Reshetnikov and the \texttt{ProveIt} development,
\emph{The Fabius Function and Rvachev's Up-Function: Exact arithmetic,
probability, Fourier analysis, and corrected sharp asymptotics}.
Primary TeX exposition, read at repository commit
\nolinkurl{2c6baf6738a5fbf3981a27141950edbe5952b03c}.
\url{https://github.com/VladimirReshetnikov/ProveIt/tree/2c6baf6738a5fbf3981a27141950edbe5952b03c/Analysis/FabiusFunction/docs/Fabius_Function_and_Rvachev_Up}.

\bibitem{webcompanion}
\texttt{ProveIt} companion exposition,
\emph{The Dyadic Web of the Fabius Function: Moment corrections,
half-base $q$-interpolation, binary telescopes, and discrete extrapolation}.
Repository snapshot dated 25 August 2026; archived under
\path{Analysis/FabiusFunction/docs/archive/research-frontiers/Fabius_Dyadic_Formulae_and_Alternative_Representations/}.
Read at the same commit as \cite{repo}.

\bibitem{arias1982}
Juan Arias de Reyna,
\emph{An infinitely differentiable function with compact support:
Definition and properties}. English translation of the paper published in
\emph{Revista de la Real Academia de Ciencias Exactas, F\'isicas y Naturales
(Madrid)} \textbf{76} (1982), 21--38.
\url{https://arxiv.org/abs/1702.05442}.

\bibitem{arias2018}
Juan Arias de Reyna,
\emph{Arithmetic of the Fabius function}.
\emph{Integers} \textbf{18} (2018), Article A51, 20 pp.
\url{https://math.colgate.edu/~integers/s51/s51.pdf}.
Preprint: \url{https://arxiv.org/abs/1702.06487}.

\bibitem{haugland}
J. K. Haugland,
\emph{Evaluating the Fabius function}.
Preprint, 2016; revised version dated 5 May 2020.
\url{https://arxiv.org/abs/1609.07999}.

\bibitem{reshetnikov2019}
Vladimir Reshetnikov,
\emph{Conjectured formula for the Fabius function}.
Mathematics Stack Exchange, question posted 5 July 2019; with subsequent
answers and discussion.
\url{https://math.stackexchange.com/questions/3283519/conjectured-formula-for-the-fabius-function}.

\bibitem{dlmfBernoulli}
NIST Digital Library of Mathematical Functions,
\S24.2, \emph{Definitions and Generating Functions: Bernoulli and Euler
Polynomials}.
\url{https://dlmf.nist.gov/24.2}.

\bibitem{dlmfq}
NIST Digital Library of Mathematical Functions,
\S17.2, \emph{$q$-Calculus}.
\url{https://dlmf.nist.gov/17.2}.

\bibitem{dlmfSinh}
NIST Digital Library of Mathematical Functions,
\S4.36, \emph{Infinite Products and Partial Fractions: Hyperbolic Functions}.
\url{https://dlmf.nist.gov/4.36}.
\end{thebibliography}
\endgroup
\end{document}
